% ===========================================================================
% Standalone implementation notes for the calo-ddpm inpainting algorithms.
% Compile:  latexmk -pdf implementation_notes.tex
% The gradient-based sections describe the CURRENT implementations
% (calo_inpaint/guided/, adapted from the CelebA-validated code).
% ===========================================================================
\documentclass[11pt]{article}
\usepackage[margin=1in]{geometry}
\usepackage{amsmath, amssymb}
\usepackage{booktabs}
\usepackage[numbers, sort&compress]{natbib}
\usepackage{hyperref}

%%%%% BEGIN MY LATEX COMMANDS%%%%%%%%%%%
\def\({\left(}
\def\){\right)}
\def\[{\left[}
\def\]{\right]}

%%%%% END MY LATEX COMMANDS%%%%%%%%%%%

\title{Implementation notes: Noise-free inpainting of sPHENIX\\
        calorimeter images with a pre-trained DDPM prior}
\date{\today}
\author{Roli Esha, Himanshu Raj}
\begin{document}
\maketitle

In these notes we spell out the details of the various in-painting algorithms based on a pre-trained un-conditional diffusion prior that we used in this work as well as how were they implemented.

\section{Setting and notation}
\label{sec:setup}

\paragraph{Data:}
From the pre-trained un-conditional diffusion model we sample 10,000 calorimeter events that are $24 \times 64$ gray-scale images. These are towers of $E_T$ in the $(\eta,\phi)$ plane, in GeV scale.  The sampling operation by the diffusion process operates in the normalized log space $x = \ln \(\max(E_T, E_{\min})\)$, $E_{\min}=10^{-3}$~GeV. The obtained sample are then mapped back by $E_T = e^{x}$ to the physical $E_T$ space. The physical log-space range used by all clamps below is $[x_{\min}, x_{\max}] = [-7,\, 4]$ ($\approx$ 0.9~MeV to 55~GeV). From \cite{torbunov2024calo} we obtained this data for both centrality 0 and centrality 4 collisions. 

\paragraph{Diffusion prior.}
Here we quickly outline the diffusion model processes. In the variance-preserving form, forward process is a sequence of successive noising process with $T = 8000$ and linear $\beta_t$ from $b/T$ to $1000\,b/T$, $b = 0.02$ (0--10\%) or $0.10$ (40--50\%); the schedule is reconstructed exactly from each checkpoint's training configuration.  With $\bar\alpha_t = \prod_{i \le t} (1 - \beta_i)$ we abbreviate
\begin{equation}
  \bar s_t = \sqrt{\bar\alpha_t}, \qquad
  \bar v_t = 1 - \bar\alpha_t ,
\end{equation}
stored in arrays of length $S{+}1$ with an identity transition padded at index~0 ($\bar s_0 = 1$, $\bar v_0 = 0$).  Sampling runs on a subsampled grid $0 = \tau_0 < \tau_1 = 1 < \dots < \tau_S = T$ ($\tau = \mathrm{round}(\mathrm{linspace}(1,T,S))$; the calorimeter study runs generation \emph{and} inpainting on the full grid, $S = T = 8000$, so the machinery below is exercised with the identity subsampling); the network is always conditioned on the \emph{original} indices $\tau_k$.  Below, one reverse step goes from level $s \equiv \tau_k$ to $s{-}1 \equiv \tau_{k-1}$, with all quantities indexed by the subsampled position.

We took the trained weights of the pre-trained DDPM~\cite{ho2020ddpm} per centrality class \cite{torbunov2024calo}. 

\paragraph{Denoiser and transport.}
The $\epsilon$-network gives the Tweedie estimate and the (jump-aware)
ancestral step used by all algorithms:
\begin{align}
  \hat x_0 &= \frac{x_s - \sqrt{\bar v_s}\,
              \epsilon_\theta(x_s, \tau_k)}{\bar s_s},
  \label{eq:x0hat}\\
  \beta_{(k)} &= 1 - \frac{\bar s_s^2}{\bar s_{s-1}^2}, \qquad
  \mathrm{ddpm\_step}(x_s, \hat x_0) =
      \frac{\bar s_{s-1}\beta_{(k)}}{\bar v_s}\,\hat x_0
    + \frac{(\bar s_s/\bar s_{s-1})\,\bar v_{s-1}}{\bar v_s}\, x_s
    + \sqrt{\tfrac{\beta_{(k)}\bar v_{s-1}}{\bar v_s}}\;\xi .
  \label{eq:ddpmstep}
\end{align}
At $s = 1$ the padding gives exact collapse,
$\mathrm{ddpm\_step} = \hat x_0$ with zero variance.  The chain is
initialized at the prior marginal $x_S \sim \mathcal{N}(0, \bar v_S
\mathbf{I})$.  The forward (re-noising) kernels used below are
$q_{s-1}(y) = \bar s_{s-1} y + \sqrt{\bar v_{s-1}}\,\xi$
(known-content branch) and the one-step jump
$x_s = (\bar s_s/\bar s_{s-1})\,x_{s-1} + \sqrt{\beta_{(k)}}\,\xi$
(RePaint time travel).

\paragraph{Problem.}
Noise-free masking: binary $m$ (1 = known/live, 0 = dead),
measurement $y = m \odot x_0^{\mathrm{true}}$, $\sigma_y = 0$; target
$p(x_0 \mid m \odot x_0 = y)$.  $A = \operatorname{diag}(m)$ is its own
pseudoinverse on its range, so no SVD machinery appears anywhere.
Posterior sampling batches $J$ samples per truth image ($J = 50$), with
$K$ images processed jointly (GPU batch $K \! \cdot \! J$); results are
stored for the dead region only, in GeV.

\section{Projection-type algorithms
         (\texttt{calo\_inpaint/inpainting/})}
\label{sec:projection}

\subsection{RePaint}
Per level $s$, repeated $U$ times ($U = 10$; a single pass at $s = 1$,
where the extra passes of the published Algorithm~1 are deterministic
recomputations):
\begin{align}
  x^{\mathrm{known}} &= q_{s-1}(y), &
  x^{\mathrm{dead}} &= \mathrm{ddpm\_step}(x_s, \hat x_0), &
  x_{s-1} &= m \, x^{\mathrm{known}} + (1{-}m)\, x^{\mathrm{dead}},
\end{align}
with the one-step forward jump between repetitions.  The known branch is
noised to level $s{-}1$ (\emph{not} $s$) so both branches coincide in
time after the step --- the off-by-one variant biases the dead-region
boundary.  Output known region equals $y$ exactly
($\bar s_0 = 1$, $\bar v_0 = 0$).

\subsection{DDNM}
Range/null-space rectification of the denoised estimate
\cite{wang2022ddnm}, exact for noiseless masking:
$\hat x_0 \leftarrow m\,y + (1{-}m)\,\hat x_0$, followed by
Eq.~(\ref{eq:ddpmstep}).  No noisy-case ($\Sigma$-scaled) corrections
are required.  Output equals the projected $\hat x_0$ at the last step;
known region exact.

\subsection{DDRM}
The spectral variational posterior of Ref.~\cite{kawar2022ddrm}
specializes (pixel basis, $\sigma_y = 0$) to
\begin{align}
  \text{dead:}\quad
  x_{s-1} &= \bar s_{s-1}\hat x_0
     + \sqrt{(1-\eta^2)\,\bar v_{s-1}}\;\hat\epsilon
     + \eta\sqrt{\bar v_{s-1}}\;\xi,\\
  \text{known:}\quad
  x_{s-1} &= \bar s_{s-1}\bigl[(1-\eta_b)\hat x_0 + \eta_b\, y\bigr]
     + \sqrt{\bar v_{s-1}}\;\xi',
\end{align}
with $\eta = 0.85$, $\eta_b = 1$ (published defaults) and matching
initialization (known pixels start at $q_S(y)$).  The known region is
additionally projected to $y$ at the final step (a no-op at
$\eta_b = 1$).  Two properties established for this noise-free limit:
(i) the dead-pixel update is the only DDIM-type transport among the
implemented algorithms, and its $\eta$ never passes through the
ancestral point ($\eta \to 1$ degenerates to re-noising $\hat x_0$);
(ii) with an \emph{exact} denoiser it under-disperses the posterior
(std ratio $\approx 0.86$ at $\eta = 0.85$; $\to 1$ as $\eta \to 0$,
$\approx 0.70$ at $\eta = 1$; independent of $S$) --- intrinsic to the
variational construction, not an implementation artifact.

\section{Gradient-guided algorithms (\texttt{calo\_inpaint/guided/})}
\label{sec:guided}

The gradient-based algorithms in production are adaptations of an
implementation developed and validated on CelebA image inpainting.
Both follow one template per reverse step; they differ only in the loss
argument and the guidance weight.  \textbf{Bookkeeping note:} these are
\emph{not} the published MCG~\cite{chung2022mcg} /
$\Pi$GDM~\cite{song2023pigdm} algorithms: both carry a RePaint-style
per-step known-region paste (absent from published $\Pi$GDM), a clamped
transport estimate, and a scalar annealed guidance weight in place of the
exact guidance coefficients.  They are accurately described as
``projection $+$ annealed consistency-gradient guidance'' and should be
cited as MCG-inspired / $\Pi$GDM-inspired variants.

\subsection{Common template}
\begin{enumerate}
  \item $\hat x_0$ from Eq.~(\ref{eq:x0hat}); clamped copy
        $\hat x_0^{c} = \mathrm{clip}(\hat x_0, x_{\min}, x_{\max})$
        used for transport.
  \item Consistency loss on the \emph{known} region, per sample
        (mean over pixel dimensions),
        \begin{equation}
          \mathcal{L} = \frac{1}{D}\bigl\lVert m \odot (\chi - y)
                        \bigr\rVert_2^2 ,
          \qquad
          g = \nabla_{x_s} \mathcal{L}
          \quad\text{(through the frozen network)},
          \label{eq:loss}
        \end{equation}
        where $\chi$ is algorithm-specific (below) and $D$ the pixel
        count.  The squared loss has a smooth gradient everywhere ---
        in particular no $0/0$ at zero residual, a failure mode of
        $\lVert\cdot\rVert_2$-based losses whose autograd gradient is
        NaN at exactly zero and at overflow.
  \item Per-sample gradient clipping,
        $g \leftarrow g \cdot \min\!\bigl(1, 1/\lVert g\rVert_2\bigr)$
        (per sample, not per batch: a global norm would shrink the
        guidance with the study batch $K\!\cdot\!J$; per-sample
        preserves the single-image behavior of the original code at any
        batch size).
  \item Guided transport on the dead region and hard paste of the known
        region at level $s{-}1$:
        \begin{align}
          x^{\mathrm{dead}} &= \mathrm{ddpm\_step}(x_s, \hat x_0^{c})
                               \;-\; \lambda_s\, g, &
          x_{s-1} &= m\, q_{s-1}(y) + (1{-}m)\, x^{\mathrm{dead}},
        \end{align}
        with $q_{0}(y) = y$; a final clean paste follows the loop.
        The annealed weight $\lambda_s$ is strong at high noise and
        vanishes as the chain sharpens.
\end{enumerate}
All gradient evaluations run in fp32 (no autocast), and every step
carries finiteness tripwires: a diverging chain raises with the step
index rather than writing NaNs to disk.

\subsection{MCG-inspired variant (\texttt{mcg2})}
Loss argument $\chi = \hat x_0^{c}$ (clamped) and
\begin{equation}
  \lambda_s = g_{\mathrm{s}} \,(1 - \bar\alpha_s),
  \qquad g_{\mathrm{s}} = 0.5 .
\end{equation}

\subsection{$\Pi$GDM-inspired variant (\texttt{pigdm2})}
Loss argument $\chi = \hat x_0$ (unclamped) and a clamped SNR-like
pseudoinverse weight,
\begin{equation}
  r_s = \min\!\left(\frac{1-\bar\alpha_s}{\bar\alpha_s + 10^{-8}},\;
        r_{\max}\right), \qquad
  \lambda_s = g_{\mathrm{s}}\,(1-\bar\alpha_s)\, r_s ,
  \qquad g_{\mathrm{s}} = 0.5,\; r_{\max} = 10 .
\end{equation}
The cap $r_{\max}$ bounds the weight where
$\bar\alpha_s \to 0$ (early chain), where the uncapped ratio reaches
$\mathcal{O}(10^4)$ on this schedule.

\subsection{Stability rationale (lessons from the discarded ports)}
\label{sec:guided:stability}

Faithful ports of published MCG/$\Pi$GDM failed on the trained
calorimeter model, and the mechanisms are worth recording:
\begin{enumerate}
  \item \emph{Amplification.}  Eq.~(\ref{eq:x0hat}) multiplies network
        error by $1/\bar s_s$, up to $\sim 150$ on this schedule; the
        guidance Jacobian $\partial \hat x_0/\partial x_s$ carries a
        second such factor.  Boundedness of exact-$\Pi$GDM-style
        guidance relies on the near-cancellation
        $(\mathbf{I} - \sqrt{\bar v_s}\,\partial\hat\epsilon/\partial
        x)/\bar s_s$, which holds for an ideal denoiser but not for the
        trained network (nor under bf16, whose $\sim 0.4\%$ error breaks
        it outright).  Without a per-step projection, the resulting
        feedback loop diverges within a few steps --- reproduced
        empirically in single precision.
  \item \emph{Norm-loss singularities.}  Losses of the form
        $\lVert r \rVert_2$ have autograd gradient $r/\lVert r\rVert$,
        which is NaN both at $r = 0$ (exact-zero residual) and at
        overflow ($\Sigma r^2 \to \infty$); the squared, per-sample-mean
        loss of Eq.~(\ref{eq:loss}) has neither.
  \item \emph{What actually stabilizes the working variant:} the
        per-step known-region paste bounds the state trajectory, the
        clamped transport estimate bounds the mean, the per-sample
        gradient clip bounds the correction, and the annealed
        $\lambda_s$ removes guidance precisely where the chain is most
        error-amplifying.  The cost is deviation from the published
        algorithms, which is why these arms are labeled
        ``-inspired'' throughout.
\end{enumerate}

\section{Common infrastructure and validation}
\label{sec:infra}

\begin{table}[h]
  \centering
  \caption{Implemented algorithms and defaults.  ``Exact $y$'': known
  region of the output equals the measurement identically.  NFEs per
  posterior sample; bwd = additional backward pass per step.}
  \label{tab:summary}
  \begin{tabular}{llllc}
    \toprule
    Algorithm & Module & Hyperparameters & Exact $y$ & NFEs \\
    \midrule
    RePaint & \texttt{inpainting/repaint} & $U{=}10$ & yes & $U(S{-}1){+}1$ \\
    DDNM    & \texttt{inpainting/ddnm}    & --- & yes & $S$ \\
    DDRM    & \texttt{inpainting/ddrm}    & $\eta{=}0.85$, $\eta_b{=}1$ & yes & $S$ \\
    MCG-inspired & \texttt{guided/mcg2}   & $g_{\mathrm{s}}{=}0.5$ & yes & $S$ (+bwd) \\
    $\Pi$GDM-inspired & \texttt{guided/pigdm2} &
        $g_{\mathrm{s}}{=}0.5$, $r_{\max}{=}10$ & yes & $S$ (+bwd) \\
    \bottomrule
  \end{tabular}
\end{table}

\paragraph{Reproducibility.}
Each study run draws from dedicated seeded generators, re-seeded per
image chunk ($\mathrm{seed} \cdot 1000003 + \mathrm{chunk\ start}$), so
runs are bitwise reproducible for fixed \texttt{--images-per-batch} and
resume exactly from the last completed chunk.  All run parameters are
recorded in a \texttt{metadata.json} next to the results.

\paragraph{Validation.}
Three independent checks back the implementations.
(i)~\emph{Schedule cross-check}: the reconstructed schedule, subsampling,
and posterior coefficients agree with the original training code of
Ref.~\cite{torbunov2024calo} to single precision.
(ii)~\emph{Analytic posterior test}: for an i.i.d.\ Gaussian prior the
optimal $\epsilon$-predictor is closed-form and the dead-pixel posterior
is known exactly; all five production algorithms reproduce it to
sampling accuracy, and all return the known region exactly.
(iii)~\emph{Adversarial boundedness}: with a random-weight (worst-case
Jacobian) network on the full $T = 8000$ schedule, the guided algorithms
remain finite and range-bounded over the entire chain --- the
configuration under which exact-$\Pi$GDM-style guidance diverges.

\bibliographystyle{unsrtnat}
\bibliography{inpainting_methods}

\end{document}
